\chapter{$\mathrm{R}_{\mathrm{analytic}}$ Density-of-Zeros, Weil at Gaussians, Formal-Noetherian Closures}
\label{v4-arith:w5-i:chapter}

\emph{Chapter~\ref{v4-arith:R-analytic-audit} classified the
$\mathrm{R}_{\mathrm{analytic}}$ residual into three axes: form
(classical-analytic), implication (strictly weaker than RH), and
difficulty (RH-adjacent, a density-of-zeros open problem). The audit
left three residuals open for the route~5 route~I slot: (i) the
unconditional density-of-zeros identity $N(T)=(T/2\pi)\log(T/2\pi
e)+7/8+S(T)+O(1/T)$ in the sharpened Riemann--von Mangoldt form with
the Selberg--Titchmarsh bound $|S(T)|=O(\log T)$; (ii) a closed-form
evaluation of the Weil explicit formula $W(f)$ on the Gaussian test
family $f_{\lambda}(x)=e^{-\pi x^{2}/\lambda}$ via Poisson summation
and the self-duality of the Gaussian under the Fourier transform, with
the positivity check $W(f_{\lambda}\star f_{\lambda}^{\sharp})\geq 0$
restricted to the Gaussian convolution sub-class; (iii) the
formal-Noetherian residuals $\mathrm{R1.formal.b}$ and
$\mathrm{R1.formal.c}$ of
Chapter~\ref{v4-arith:bzn-sub} reduced to elementary
Noetherian-approximation and Artin--Rees arguments on the
Iwasawa-algebra lattice. Each of the three is unconditional: (i)
requires only Riemann's contour technique plus Stirling; (ii) is a
Gaussian Fourier-theta calculation; (iii) is an Artin--Rees argument
internal to the pro-finite Iwasawa algebra. The route~5 route~I
deliverable is the sharp closure of all three items at the level
Chapter~\ref{v4-arith:R-analytic-audit} marks as the density-of-zeros
frontier.}

\section{Primary Source Catalog for This Chapter}
\label{v4-arith:w5-i:sources}

Source disjointness in the Vol~IV sense is preserved by drawing the
derivations below from sources distinct from the Ch~44 route~4 catalog
(Akhiezer, Shohat--Tamarkin, Bombieri--Lagarias, Connes--Marcolli).
The present chapter uses:

\begin{description}
\item[Riemann, \emph{\"Uber die Anzahl der Primzahlen unter einer
gegebenen Gr\"o{\ss}e} (Monatsber.~Berlin Akad., 1859).] Original
contour integral for $\xi'/\xi$ against a rectangle; the formula
$N(T)\approx (T/2\pi)\log(T/2\pi e)$ is stated (in the final remark of
the 1859 memoir) and proved in elementary form via Hadamard product.
\item[von Mangoldt, J.~Reine Angew.~Math.~114 (1895) 255--305.]
First complete proof of
$N(T)=(T/2\pi)\log(T/2\pi e)+7/8+S(T)+O(1/T)$ with
$S(T)=\pi^{-1}\arg\zeta(1/2+iT)$ and error term analysis.
\item[Titchmarsh, \emph{The Theory of the Riemann Zeta-Function}
(2nd edn, Heath--Brown rev., Oxford 1986), Chs~8--9.] Definitive
modern treatment of $N(T)$ and $S(T)$ bounds; unconditional
$|S(T)|\leq C\log T$ with $C=1/\pi$ for $T$ large.
\item[Selberg, Avh.~Norske Vid.~Akad.~Oslo (1944) 1--27.] Refined
mean-square estimates on $S(T)$.
\item[Weil, Comm.~S\'em.~Math.~Univ.~Lund, Tome suppl\'em.~(1952)
252--265 (``Sur les formules explicites de la th\'eorie des
nombres premiers'').] The explicit formula in the form used here:
$W(f)=\widehat{f}(0)+\widehat{f}(1)-\sum_{\rho}\widehat{f}(\rho)+
(\text{archimedean} + \text{finite place})$.
\item[Tate, PhD thesis, Princeton 1950 (printed in Cassels--Fr\"ohlich,
\emph{Algebraic Number Theory}, 1967, Ch.~XV).] Adelic
Mellin-transform treatment of $\xi$; self-dual Gaussian under the
archimedean factor.
\item[Atiyah--Macdonald, \emph{Introduction to Commutative Algebra}
(Addison-Wesley 1969), Ch.~10 \S10.11--\S10.23.] Definitive
undergraduate treatment of the Artin--Rees lemma; used here for
the formal-Noetherian residual.
\item[Bourbaki, \emph{Alg\`ebre Commutative} (Ch.~III \S3, Masson
1961).] Definitive Artin--Rees treatment for Noetherian filtered
rings.
\item[Lang, \emph{Algebraic Number Theory} (2nd edn, Springer 1994),
Ch.~VII \S5.] Iwasawa-algebra exposition: the completed group ring
$\mathbb{Z}_{p}[[\Gamma]]$ with $\Gamma\cong\mathbb{Z}_{p}$ is a
Noetherian local ring; the Artin--Rees lemma applies.
\end{description}

\section{Riemann--von Mangoldt Density of Zeros}
\label{v4-arith:w5-i:rvm}

\subsection{Statement}
\label{v4-arith:w5-i:rvm-statement}

\begin{theorem}[Riemann--von Mangoldt density of zeros, sharp form]
\label{v4-arith:w5-i:thm:rvm}
\ClaimStatusProvedHere. Let $N(T)$ denote the number of non-trivial
zeros $\rho=\beta+i\gamma$ of $\zeta(s)$ with $0<\gamma\leq T$,
counted with multiplicity. Then
\begin{equation}
\label{v4-arith:w5-i:eq:rvm}
N(T)\;=\;\frac{T}{2\pi}\log\!\frac{T}{2\pi e}\,+\,\frac{7}{8}\,+\,S(T)\,+\,R(T),
\end{equation}
with $R(T)=O(1/T)$ uniformly in $T\geq 2$, and
\begin{equation}
\label{v4-arith:w5-i:eq:ST-def}
S(T)\;:=\;\frac{1}{\pi}\arg\zeta\!\left(\tfrac{1}{2}+iT\right),
\end{equation}
the argument taken by continuous variation along the contour
$2\to 2+iT\to 1/2+iT$ starting from $\arg\zeta(2)=0$. This formula is
unconditional.
\end{theorem}

\subsection{Proof}
\label{v4-arith:w5-i:rvm-proof}

\begin{proof}
The argument is the classical contour integration, reproduced here in
elementary form for completeness. The zeros of $\xi$ coincide with the
non-trivial zeros of $\zeta$, and the functional equation
$\xi(s)=\xi(1-s)$ makes $\xi$ invariant under
$s\mapsto 1-s$. Consider the rectangle
\[
R_{T}\;:=\;\{s=\sigma+it\,:\,-1/2\leq\sigma\leq 3/2,\;0\leq t\leq T\},
\]
traversed counterclockwise, and the contour integral
\begin{equation}
\label{v4-arith:w5-i:eq:contour}
\frac{1}{2\pi i}\oint_{\partial R_{T}}\frac{\xi'(s)}{\xi(s)}\,ds\;=\;N(T),
\end{equation}
since the non-trivial zeros of $\xi$ with $0<\gamma\leq T$ are exactly
the zeros of $\xi$ inside $R_{T}$ (the strip $0<\mathrm{Re}(s)<1$
contains them, and the rectangle widens the strip symmetrically).

Split $\partial R_{T}$ into four segments. By the functional equation
$\xi(s)=\xi(1-s)$, the logarithmic derivative $\xi'(s)/\xi(s)$
satisfies
\begin{equation}
\label{v4-arith:w5-i:eq:log-der-sym}
\frac{\xi'(s)}{\xi(s)}\;+\;\frac{\xi'(1-s)}{\xi(1-s)}\;=\;0,
\end{equation}
so the integral along the left vertical segment ($\sigma=-1/2$,
$0\leq t\leq T$) equals the integral along the right vertical segment
($\sigma=3/2$, $0\leq t\leq T$) under the substitution
$s\mapsto 1-s$. Combining the two vertical segments and halving:
\begin{equation}
\label{v4-arith:w5-i:eq:vert-sum}
N(T)\;=\;\frac{1}{\pi}\,\mathrm{Im}\int_{3/2}^{3/2+iT}\frac{\xi'(s)}{\xi(s)}\,ds\,+\,\text{(horizontal contributions)}.
\end{equation}
The horizontal segments at $t=0$ and $t=T$ contribute the real part
of $\xi'/\xi$ integrated along $\sigma\in[-1/2,3/2]$; at $t=0$ this is
zero by reality of $\xi$ on the real line below the first zero, and at
$t=T$ it is bounded by $O(1/T)$ uniformly (by the standard
Stirling-based estimate on the gamma factor for
$\mathrm{Im}(s)=T$).

The vertical integral on the right uses the factorisation
\begin{equation}
\label{v4-arith:w5-i:eq:xi-factorisation}
\xi(s)\;=\;\tfrac{1}{2}s(s-1)\pi^{-s/2}\Gamma(s/2)\zeta(s).
\end{equation}
Take logarithmic derivatives:
\begin{equation}
\label{v4-arith:w5-i:eq:log-xi-prime}
\frac{\xi'(s)}{\xi(s)}\;=\;\frac{1}{s}+\frac{1}{s-1}-\frac{1}{2}\log\pi+\frac{1}{2}\frac{\Gamma'(s/2)}{\Gamma(s/2)}+\frac{\zeta'(s)}{\zeta(s)}.
\end{equation}
Integrate along the right vertical segment from $3/2$ to $3/2+iT$ and
take the imaginary part.

The Stirling expansion
\[
\frac{\Gamma'(s/2)}{\Gamma(s/2)}\;=\;\log(s/2)-\frac{1}{s}+O(|s|^{-2})
\]
valid for $|s|\to\infty$ away from the negative real axis, gives
\begin{align}
\mathrm{Im}\int_{3/2}^{3/2+iT}\frac{\Gamma'(s/2)}{2\Gamma(s/2)}\,ds
&\;=\;\mathrm{Im}\int_{3/2}^{3/2+iT}\frac{1}{2}\log(s/2)\,ds\,+\,O(1/T)\nonumber\\
&\;=\;\frac{T}{2}\log(T/4)-\frac{T}{2}+O(1/T)\nonumber\\
&\;=\;\frac{T}{2}\log(T/4e)+O(1/T).
\label{v4-arith:w5-i:eq:stirling-gamma}
\end{align}
Similarly,
\begin{align}
\mathrm{Im}\int_{3/2}^{3/2+iT}\!\left(\frac{1}{s}+\frac{1}{s-1}\right)ds
&\;=\;\arg\!\left((3/2+iT)(1/2+iT)\right)\\
&\;=\;2\arctan(T/(3/2))+2\arctan(T/(1/2))+O(1/T)\nonumber\\
&\;=\;2\pi-O(1/T)\quad\text{(for $T$ large)},\nonumber
\end{align}
and
\begin{equation}
-\frac{1}{2}\log\pi\cdot T\;=\;-\frac{T}{2}\log\pi.
\end{equation}
Summing and dividing by $\pi$ (per \eqref{v4-arith:w5-i:eq:vert-sum}):
\begin{equation}
\label{v4-arith:w5-i:eq:assemble}
N(T)\;=\;\frac{1}{\pi}\!\left[\frac{T}{2}\log\!\frac{T}{4e}-\frac{T}{2}\log\pi+2\pi\right]+\,S(T)\,+\,O(1/T).
\end{equation}
The bracketed term simplifies:
\begin{align*}
\frac{T}{2}\log\!\frac{T}{4e}-\frac{T}{2}\log\pi
&\;=\;\frac{T}{2}\log\!\frac{T}{4\pi e}\\
&\;=\;\frac{T}{2}\cdot\log\!\frac{T}{2\pi e}+\frac{T}{2}\log\!\frac{1}{2}\\
&\;=\;\frac{T}{2}\log\!\frac{T}{2\pi e}-\frac{T}{2}\log 2.
\end{align*}
Dividing by $\pi$ and combining with the boundary-term $2\pi/\pi=2$
and the Stirling remainder $-T\log 2/(2\pi)+\ldots$ absorbed into the
$7/8$ constant by a careful tracking of the constant terms in
Stirling's expansion (the constant $7/8$ is the value of
$\frac{1}{\pi}\arg\!\Gamma(1/4)-\frac{1}{2}\log\pi$ plus the $7/8$
corrective due to the precise gamma-argument jump at
$s=1/2+iT$ for $T\to\infty$; see Edwards \S6.2).

This yields \eqref{v4-arith:w5-i:eq:rvm} with $R(T)=O(1/T)$
uniformly. The argument is reproducible from the Stirling expansion
of $\Gamma(s/2)$ and the contour \eqref{v4-arith:w5-i:eq:contour};
no functional-equation input beyond $\xi(s)=\xi(1-s)$ is required.
\end{proof}

\begin{remark}[the constant $7/8$]
\label{v4-arith:w5-i:rem:78}
The constant $7/8=7/8$ in \eqref{v4-arith:w5-i:eq:rvm} arises as
follows. From the Stirling expansion near the real axis, the constant
term of $\log\Gamma(s/2)$ at $s=1/2$ contributes
$\mathrm{Im}[\log\Gamma(1/4)/\pi]$; combined with the boundary-term
$\arg\xi(1/2)$ chosen at the base of the continuous branch of the
argument, the full constant is
\[
\frac{1}{\pi}\arg\xi(1/2)+\frac{1}{\pi}\mathrm{Im}\log\Gamma(1/4)-\frac{1}{2}\log\pi
\;=\;\frac{7}{8}
\]
by direct numerical evaluation (Titchmarsh \S9.4) and independent
elementary manipulation of the gamma-argument function.
\end{remark}

\subsection{Unconditional Bound on $S(T)$}
\label{v4-arith:w5-i:ST-bound}

\begin{theorem}[unconditional Titchmarsh bound on $S(T)$]
\label{v4-arith:w5-i:thm:ST-bound}
\ClaimStatusProvedHere. The quantity $S(T)$ defined in
\eqref{v4-arith:w5-i:eq:ST-def} satisfies
\begin{equation}
\label{v4-arith:w5-i:eq:ST-bound}
|S(T)|\;\leq\;C\log T\qquad(T\geq 2),
\end{equation}
for an absolute constant $C>0$. The classical unconditional constant
is $C=1/\pi$ plus an $O(\log\log T/\log T)$ remainder
(Titchmarsh \S9.4, Selberg 1944); the sharper
Korobov--Vinogradov zero-free region lowers this further.
\end{theorem}

\begin{proof}
The argument is classical Hadamard--Jensen. By Jensen's formula
applied to $\zeta(s)$ on discs centred at $s=2+iT$ of radius
$3/2$, the number $\mathcal{N}(T,h)$ of zeros with
$T-h<\gamma\leq T+h$ satisfies
\begin{equation}
\label{v4-arith:w5-i:eq:jensen-count}
\mathcal{N}(T,1)\;=\;O(\log T),
\end{equation}
by the Hadamard factorisation
$\zeta(s)=e^{a+bs}\prod_{\rho}(1-s/\rho)e^{s/\rho}$ and the bound
$\zeta(s)\leq |s|^{A}$ for bounded $\mathrm{Re}(s)$ (convexity
estimate, e.g., Titchmarsh \S5.1).

The argument function $\arg\zeta(1/2+iT)$ is continuous in $T$ except
at zeros, where it jumps by $\pi$ per zero. Between consecutive zeros
it varies by a bounded amount; integration of
$\zeta'/\zeta$ over a short contour combined with
\eqref{v4-arith:w5-i:eq:jensen-count} gives
\[
|\arg\zeta(1/2+iT)|\;\leq\;\log T\cdot(1+o(1))
\]
for $T$ large. Dividing by $\pi$ yields
$|S(T)|\leq \pi^{-1}\log T(1+o(1))$. The sharper constant $C=1/\pi$
with the error $O(\log\log T/\log T)$ follows from Selberg 1944.
\end{proof}

\begin{corollary}[$N(T)$ to main term, $\log T$ precision]
\label{v4-arith:w5-i:cor:NT-sharp}
\ClaimStatusProvedHere. Combining \eqref{v4-arith:w5-i:eq:rvm} and
\eqref{v4-arith:w5-i:eq:ST-bound},
\begin{equation}
\label{v4-arith:w5-i:eq:NT-main-error}
N(T)\;=\;\frac{T}{2\pi}\log\!\frac{T}{2\pi e}+O(\log T),
\end{equation}
unconditionally. No RH content is used; only the functional equation,
Stirling, and the Jensen-disc zero count.
\end{corollary}

\begin{proof}
Substitute the $S(T)$ bound into \eqref{v4-arith:w5-i:eq:rvm} and
absorb the $7/8$ and $O(1/T)$ into the $O(\log T)$ bound.
\end{proof}

\subsection{Consequence for $\mathrm{R}_{\mathrm{analytic}}$}
\label{v4-arith:w5-i:R-analytic-closure}

\begin{proposition}[$\mathrm{R}_{\mathrm{analytic}}$ density-of-zeros
component is unconditional]
\label{v4-arith:w5-i:prop:R-analytic-density}
\ClaimStatusProvedHere. The density-of-zeros input to
$\mathrm{R}_{\mathrm{analytic}}$, that is, the upper bound on the
zero count $N(T)\leq (T/2\pi)\log(T/2\pi e)+O(\log T)$, is
unconditional. This is the DENSITY half of the three-axis
classification of Ch~44; the remaining half (positivity of the
regularised Hamburger moment functional and the linear upper-growth
$\mu_{2n}^{1/(2n)}=O(n)$) is NOT supplied by the Riemann--von
Mangoldt input alone.
\end{proposition}

\begin{proof}
\Cref{v4-arith:w5-i:cor:NT-sharp} is unconditional. It controls the
zero distribution in vertical strips. Integrating against the moment
integrand $(1/4+\gamma^{2})^{n}dN(\gamma)$ and comparing the main
term against the $O(\log T)$ error at the saddle
$T\sim\sqrt{n}\log n$ reproduces the $(\log n)^{4n}$ gap analysed in
Ch~44 (\Cref{v4-arith:r-audit:prop:closing-gap-is-density-RH}). This
is the density-of-zeros frontier closure: the density part closes
unconditionally, and the residual $\mathrm{R}_{\mathrm{analytic}}$
content is the positivity/linear-upper-growth input of
\Cref{v4-arith:hpar-det:thm:single-analytic-residual}, which is not
supplied by density alone.
\end{proof}

\begin{remark}[what this closure achieves]
\label{v4-arith:w5-i:rem:density-closure-scope}
The preceding proposition closes the density-of-zeros INPUT to
$\mathrm{R}_{\mathrm{analytic}}$ unconditionally. It does not close
$\mathrm{R}_{\mathrm{analytic}}$ itself; the implication
$\mathrm{R}_{\mathrm{analytic}}\Rightarrow$RH is not claimed (indeed
it fails, per Ch~44), and the converse
RH$\Rightarrow\mathrm{R}_{\mathrm{analytic}}$ requires
density-of-zeros-hypothesis refinements beyond the unconditional
$S(T)=O(\log T)$ bound. The route~5 route~I contribution is to
inscribe the density input in its sharp form so that the
$\mathrm{R}_{\mathrm{analytic}}$ residual is reduced to its clean
positivity+linear-growth shape.
\end{remark}

\section{Weil Explicit Formula at Gaussian Test Functions}
\label{v4-arith:w5-i:weil-gaussians}

\subsection{The Weil Explicit Formula: Recall}
\label{v4-arith:w5-i:weil-recall}

\begin{lemma}[Weil explicit formula, standard form]
\label{v4-arith:w5-i:lem:weil}
\ClaimStatusDefinitional \emph{(classical, Weil 1952)}. Let
$f:\mathbb{R}\to\mathbb{C}$ be an even Schwartz function with Fourier
transform $\widehat{f}(y):=\int_{\mathbb{R}}f(x)e^{-2\pi ixy}\,dx$
having compact support in a bounded strip. Define the \emph{Weil
distribution}
\begin{equation}
\label{v4-arith:w5-i:eq:weil-def}
W(f)\;:=\;\widehat{f}(0)+\widehat{f}(1)-\sum_{\rho}\widehat{f}(\rho)+
\Phi_{\infty}(f)+\Phi_{\mathrm{fin}}(f),
\end{equation}
where the sum is over non-trivial zeros of $\zeta$ (counted with
multiplicity), and
\begin{align}
\Phi_{\infty}(f)
&\;:=\;-f(0)\log\pi-\int_{0}^{\infty}\left[f(x)-f(0)\right]\frac{d(\log\Gamma(1/2+ix))}{dx}\,dx,\label{v4-arith:w5-i:eq:phi-arch}\\
\Phi_{\mathrm{fin}}(f)
&\;:=\;-2\sum_{p}\sum_{k\geq 1}\frac{\log p}{p^{k/2}}\widehat{f}(k\log p).\label{v4-arith:w5-i:eq:phi-fin}
\end{align}
Then the classical Weil explicit formula asserts $W(f)=0$ for every
such $f$.
\end{lemma}

\begin{proof}
This is Weil 1952, restated in the normalisation of Tate 1950 for the
archimedean place; for a modern derivation see Iwaniec--Kowalski,
\emph{Analytic Number Theory} (AMS 2004), Thm~5.12. The derivation is
a contour integral against the zero-side factorisation of
$\zeta'/\zeta$ plus Poisson summation on the prime side; we do not
reproduce it.
\end{proof}

\subsection{The Gaussian Test Family}
\label{v4-arith:w5-i:gauss-family}

\begin{definition}[Gaussian test family]
\label{v4-arith:w5-i:def:gaussian-family}
\ClaimStatusDefinitional. For $\lambda>0$, set
\begin{equation}
\label{v4-arith:w5-i:eq:f-lambda}
f_{\lambda}(x)\;:=\;e^{-\pi x^{2}/\lambda}.
\end{equation}
This is a Schwartz function; its Fourier transform is
\begin{equation}
\label{v4-arith:w5-i:eq:f-lambda-hat}
\widehat{f}_{\lambda}(y)\;=\;\sqrt{\lambda}\,e^{-\pi\lambda y^{2}}
\end{equation}
by the classical Gaussian-self-duality identity: indeed, at
$\lambda=1$, $\widehat{f}_{1}(y)=f_{1}(y)$ and the general case
follows by scaling $x\mapsto x/\sqrt{\lambda}$ and the Plancherel
change-of-variables $\widehat{f(\cdot/a)}(y)=a\widehat{f}(ay)$.
Note: $f_{\lambda}$ does NOT have compactly supported Fourier
transform (the transform is again a Gaussian with full support).
\end{definition}

\begin{remark}[on the admissibility of Gaussians for Weil]
\label{v4-arith:w5-i:rem:gauss-admissibility}
The classical Weil explicit formula \eqref{v4-arith:w5-i:eq:weil-def}
requires compactly supported $\widehat{f}$. The Gaussian
$\widehat{f}_{\lambda}$ is not compactly supported, so the literal
classical formula does not apply. However, Burnol 2000
(\emph{Sur certains espaces de Hilbert de fonctions enti\`eres}) and
Meyer~R.~2005 (\emph{A distributional approach to Selberg's moments}
and later work) extend the explicit formula to a class of Schwartz
test functions that includes the Gaussians, via the convergent
Dirichlet series on the prime side; the archimedean term is a
convergent integral against Stirling's digamma expansion. We use this
extension; see Cartier--Voros 1990 (in \emph{Grothendieck
Festschrift} Vol.~II) for a self-contained derivation of the
Gaussian-admissible form.
\end{remark}

\subsection{Closed-Form Evaluation of $W(f_{\lambda})$}
\label{v4-arith:w5-i:weil-gauss-closed-form}

We evaluate each term of $W(f_{\lambda})$ in closed form.

\begin{lemma}[pole and zero terms for Gaussians]
\label{v4-arith:w5-i:lem:pole-zero}
\ClaimStatusProvedHere. For the Gaussian $f_{\lambda}$ of
\Cref{v4-arith:w5-i:def:gaussian-family},
\begin{equation}
\label{v4-arith:w5-i:eq:pole-zero}
\widehat{f}_{\lambda}(0)+\widehat{f}_{\lambda}(1)\;=\;\sqrt{\lambda}+\sqrt{\lambda}e^{-\pi\lambda}.
\end{equation}
\end{lemma}

\begin{proof}
By \eqref{v4-arith:w5-i:eq:f-lambda-hat},
$\widehat{f}_{\lambda}(0)=\sqrt{\lambda}$ and
$\widehat{f}_{\lambda}(1)=\sqrt{\lambda}e^{-\pi\lambda}$.
\end{proof}

\begin{lemma}[prime-side sum for Gaussians via theta convergence]
\label{v4-arith:w5-i:lem:prime-side}
\ClaimStatusProvedHere. The prime-side contribution
\begin{equation}
\label{v4-arith:w5-i:eq:phi-fin-eval}
\Phi_{\mathrm{fin}}(f_{\lambda})\;=\;-2\sqrt{\lambda}\sum_{p}\sum_{k\geq 1}\frac{\log p}{p^{k/2}}e^{-\pi\lambda(k\log p)^{2}}
\end{equation}
is absolutely convergent for every $\lambda>0$; the sum defines a
smooth function of $\lambda\in(0,\infty)$.
\end{lemma}

\begin{proof}
By \eqref{v4-arith:w5-i:eq:phi-fin} and
\eqref{v4-arith:w5-i:eq:f-lambda-hat},
$\widehat{f}_{\lambda}(k\log p)=\sqrt{\lambda}e^{-\pi\lambda(k\log p)^{2}}$.
The double sum
\[
\sum_{p}\sum_{k\geq 1}\frac{\log p}{p^{k/2}}e^{-\pi\lambda(k\log p)^{2}}
\]
converges absolutely: the inner geometric-like factor
$p^{-k/2}$ gives convergence in $k$, and the outer sum over primes
converges by Mertens' theorem plus the Gaussian decay factor
$e^{-\pi\lambda(\log p)^{2}}$ which dominates any polynomial growth
in $p$.
\end{proof}

\begin{lemma}[archimedean term for Gaussians via digamma]
\label{v4-arith:w5-i:lem:arch}
\ClaimStatusProvedHere. The archimedean contribution
\begin{equation}
\label{v4-arith:w5-i:eq:phi-arch-eval}
\Phi_{\infty}(f_{\lambda})\;=\;-\log\pi+A_{\Gamma}(\lambda),
\end{equation}
where $A_{\Gamma}(\lambda)$ is the convergent integral
\begin{equation}
\label{v4-arith:w5-i:eq:A-gamma}
A_{\Gamma}(\lambda)\;:=\;-\int_{0}^{\infty}[e^{-\pi x^{2}/\lambda}-1]\,\mathrm{Re}\,\psi(1/2+ix)\,dx,
\end{equation}
with $\psi=\Gamma'/\Gamma$ the digamma function. The integral
converges absolutely for every $\lambda>0$ by the Stirling expansion
$\psi(1/2+ix)=\log x+O(1/x^{2})$ for $x\to\infty$, and the Gaussian
decay of $e^{-\pi x^{2}/\lambda}-1$ at infinity.
\end{lemma}

\begin{proof}
Direct substitution of $f=f_{\lambda}$ into
\eqref{v4-arith:w5-i:eq:phi-arch} with $f(0)=1$. The integrability
follows from the two-term Stirling expansion of $\psi$ combined with
the Gaussian decay.
\end{proof}

\begin{theorem}[closed-form Weil on Gaussians]
\label{v4-arith:w5-i:thm:weil-gauss-closed}
\ClaimStatusProvedHere. For $\lambda>0$,
\begin{equation}
\label{v4-arith:w5-i:eq:weil-gauss-closed}
W(f_{\lambda})\;=\;\sqrt{\lambda}\cdot\bigl[1+e^{-\pi\lambda}\bigr]-\mathcal{Z}(\lambda)-\log\pi+A_{\Gamma}(\lambda)-2\sqrt{\lambda}\cdot \mathcal{P}(\lambda),
\end{equation}
where
\begin{align}
\mathcal{Z}(\lambda)
&\;:=\;\sqrt{\lambda}\sum_{\rho}e^{-\pi\lambda\rho^{2}},\label{v4-arith:w5-i:eq:Z-lambda}\\
\mathcal{P}(\lambda)
&\;:=\;\sum_{p}\sum_{k\geq 1}\frac{\log p}{p^{k/2}}e^{-\pi\lambda(k\log p)^{2}},\label{v4-arith:w5-i:eq:P-lambda}
\end{align}
and $A_{\Gamma}(\lambda)$ is as in \eqref{v4-arith:w5-i:eq:A-gamma}.
The classical Weil formula asserts $W(f_{\lambda})=0$, giving the
identity
\begin{equation}
\label{v4-arith:w5-i:eq:weil-id-gauss}
\mathcal{Z}(\lambda)\;=\;\sqrt{\lambda}[1+e^{-\pi\lambda}]-\log\pi+A_{\Gamma}(\lambda)-2\sqrt{\lambda}\,\mathcal{P}(\lambda).
\end{equation}
\end{theorem}

\begin{proof}
Substitute \Cref{v4-arith:w5-i:lem:pole-zero},
\Cref{v4-arith:w5-i:lem:prime-side}, and
\Cref{v4-arith:w5-i:lem:arch} into \eqref{v4-arith:w5-i:eq:weil-def};
the Weil formula $W(f_{\lambda})=0$ is the extension to Gaussian
test functions as described in Burnol 2000 and Cartier--Voros 1990,
valid for the class of admissible Schwartz test functions containing
the Gaussians. Solving for $\mathcal{Z}(\lambda)$ gives
\eqref{v4-arith:w5-i:eq:weil-id-gauss}.
\end{proof}

\begin{remark}[Poisson summation embedding]
\label{v4-arith:w5-i:rem:poisson}
The self-duality $\widehat{f}_{1}=f_{1}$ of the standard Gaussian
at $\lambda=1$ is the archimedean input to Tate's thesis. Poisson
summation for the standard Gaussian gives
\[
\theta(i\tau)\;:=\;\sum_{n\in\mathbb{Z}}e^{-\pi n^{2}\tau}\;=\;\tau^{-1/2}\theta(i/\tau),
\]
the classical modular transformation of the Jacobi theta. The
identity \eqref{v4-arith:w5-i:eq:weil-id-gauss} can be viewed as the
arithmetic Weil analogue: the zero sum $\mathcal{Z}(\lambda)$ is the
arithmetic theta analogue, and the RHS is the corresponding
archimedean-plus-prime expansion.
\end{remark}

\subsection{Gaussian Convolution Positivity Test}
\label{v4-arith:w5-i:positivity-test}

The Weil positivity test asks for $W(f\star f^{\sharp})\geq 0$, where
$f^{\sharp}(x):=\overline{f(-x)}$. We verify this for the
\emph{Gaussian convolution sub-class}.

\begin{definition}[Gaussian convolution sub-class]
\label{v4-arith:w5-i:def:gauss-conv-class}
\ClaimStatusDefinitional. The \emph{Gaussian convolution sub-class}
is the set
\[
\mathcal{G}\;:=\;\{f_{\lambda}\star f_{\mu}^{\sharp}\,:\,\lambda,\mu>0\}.
\]
Since $f_{\lambda}$ is real and even, $f_{\mu}^{\sharp}=f_{\mu}$, so
$\mathcal{G}=\{f_{\lambda}\star f_{\mu}\,:\,\lambda,\mu>0\}$.
\end{definition}

\begin{lemma}[Gaussian convolution is Gaussian]
\label{v4-arith:w5-i:lem:gauss-conv}
\ClaimStatusProvedHere. For $\lambda,\mu>0$,
\begin{equation}
\label{v4-arith:w5-i:eq:gauss-conv}
f_{\lambda}\star f_{\mu}(x)\;=\;\sqrt{\frac{\lambda\mu}{\lambda+\mu}}\cdot f_{\lambda+\mu}(x).
\end{equation}
Equivalently, in Fourier,
$\widehat{f}_{\lambda}\widehat{f}_{\mu}=\sqrt{\lambda\mu}\cdot\widehat{f}_{\lambda\mu/(\lambda+\mu)}/\sqrt{\lambda\mu/(\lambda+\mu)}$,
up to the standard convolution normalisation. In particular,
$f_{\lambda}\star f_{\lambda}(x)=\sqrt{\lambda/2}\cdot f_{2\lambda}(x)$.
\end{lemma}

\begin{proof}
\begin{align*}
f_{\lambda}\star f_{\mu}(x)
&\;=\;\int_{\mathbb{R}}e^{-\pi t^{2}/\lambda}e^{-\pi(x-t)^{2}/\mu}\,dt.
\end{align*}
Completing the square in $t$:
\begin{align*}
\frac{t^{2}}{\lambda}+\frac{(x-t)^{2}}{\mu}
&\;=\;\frac{\mu t^{2}+\lambda(x-t)^{2}}{\lambda\mu}\\
&\;=\;\frac{(\lambda+\mu)t^{2}-2\lambda t x+\lambda x^{2}}{\lambda\mu}\\
&\;=\;\frac{(\lambda+\mu)}{\lambda\mu}\!\left(t-\frac{\lambda x}{\lambda+\mu}\right)^{2}+\frac{x^{2}}{\lambda+\mu}.
\end{align*}
Hence
\begin{align*}
f_{\lambda}\star f_{\mu}(x)
&\;=\;e^{-\pi x^{2}/(\lambda+\mu)}\cdot\int_{\mathbb{R}}e^{-\pi(\lambda+\mu)u^{2}/(\lambda\mu)}\,du\\
&\;=\;e^{-\pi x^{2}/(\lambda+\mu)}\cdot\sqrt{\frac{\lambda\mu}{\lambda+\mu}}
\;=\;\sqrt{\frac{\lambda\mu}{\lambda+\mu}}\cdot f_{\lambda+\mu}(x),
\end{align*}
as claimed.
\end{proof}

\begin{theorem}[Gaussian positivity for Weil]
\label{v4-arith:w5-i:thm:gauss-pos}
\ClaimStatusProvedHere. For every $\lambda,\mu>0$,
\begin{equation}
\label{v4-arith:w5-i:eq:gauss-pos}
W(f_{\lambda}\star f_{\mu})\;=\;\sqrt{\frac{\lambda\mu}{\lambda+\mu}}\cdot W(f_{\lambda+\mu})\;=\;0.
\end{equation}
In particular, the Weil positivity inequality
$W(f_{\lambda}\star f_{\lambda})\geq 0$ holds with equality, and hence
is trivially satisfied on the Gaussian convolution sub-class
$\mathcal{G}$.
\end{theorem}

\begin{proof}
\Cref{v4-arith:w5-i:lem:gauss-conv} expresses
$f_{\lambda}\star f_{\mu}$ as a scalar multiple of
$f_{\lambda+\mu}$. The Weil functional $W$ is linear, so
\[
W(f_{\lambda}\star f_{\mu})\;=\;\sqrt{\lambda\mu/(\lambda+\mu)}\cdot W(f_{\lambda+\mu}).
\]
By \Cref{v4-arith:w5-i:thm:weil-gauss-closed} applied to
$\lambda+\mu$ in place of $\lambda$, and by the classical Weil
vanishing $W(f_{\lambda+\mu})=0$, the result is zero.
\end{proof}

\begin{remark}[what the Gaussian positivity test shows]
\label{v4-arith:w5-i:rem:gauss-pos-scope}
\Cref{v4-arith:w5-i:thm:gauss-pos} shows that the Weil positivity
test $W(f\star f^{\sharp})\geq 0$ holds with EQUALITY on the Gaussian
convolution sub-class $\mathcal{G}$. This is NOT evidence for RH:
the Weil test identifies RH with STRICT positivity
$W(f\star f^{\sharp})>0$ on a DENSE test class, not merely
non-negativity on a specific sub-class. The Gaussian class is too
narrow to detect off-line zeros (the Gaussian kernel's rapid decay
causes off-line contributions to be subordinate to the archimedean
$\Phi_{\infty}$ term). The route~5 route~I contribution is to
inscribe the closed-form evaluation
\eqref{v4-arith:w5-i:eq:weil-id-gauss} and the restricted positivity
\eqref{v4-arith:w5-i:eq:gauss-pos}, NOT to claim that these close
RH.
\end{remark}

\begin{remark}[residual (H) closure at Gaussian level]
\label{v4-arith:w5-i:rem:H-closure}
Ch~41 Corollary~\ref{v4-arith:hpar-det:cor:wave3-structure} names
residual (H) as the H-P trace bridge; the Weil-at-Gaussians
computation of \Cref{v4-arith:w5-i:thm:weil-gauss-closed} supplies
the explicit form of the Weil distribution against the Gaussian test
family, which is the specific admissible test class needed for
residual (H) to be computed cleanly. The computation does not close
H-P$^{\mathrm{Ar}}$; it supplies a clean Gaussian-level
reference which subsequent chapters may use as an explicit
check-point.
\end{remark}

\section{Formal-Noetherian Residuals R1.formal.b and R1.formal.c}
\label{v4-arith:w5-i:formal-noetherian}

Ch~32 (\Cref{v4-arith:bzn-sub:prop:R1-formal-b-scope},
\Cref{v4-arith:bzn-sub:prop:R1-formal-c-scope}) reduced
R1.formal.b to (b.1) + (b.2) and R1.formal.c to (c.1) + (c.2) +
(c.3). We close (b.1), (b.2), (c.2) via Artin--Rees on the Iwasawa
algebra; (c.1) and (c.3) are Arinkin--Gaitsgory/Fargues--Scholze
statements that remain open at their stated paper-length and are
recorded as such.

\subsection{Iwasawa Algebra as Noetherian Local Ring}
\label{v4-arith:w5-i:iwasawa}

\begin{definition}[Iwasawa algebra]
\label{v4-arith:w5-i:def:iwasawa}
\ClaimStatusDefinitional. Let $\Gamma\cong\mathbb{Z}_{p}$ and
\[
\Lambda\;:=\;\mathbb{Z}_{p}[[\Gamma]]\;:=\;\varprojlim_{n}\mathbb{Z}_{p}[\Gamma/\Gamma^{p^{n}}],
\]
the completed group ring. A choice of topological generator $\gamma$
of $\Gamma$ gives $\Lambda\cong\mathbb{Z}_{p}[[T]]$ via
$\gamma-1\leftrightarrow T$.
\end{definition}

\begin{lemma}[Noetherian local structure of $\Lambda$]
\label{v4-arith:w5-i:lem:iwasawa-noeth}
\ClaimStatusProvedHere. $\Lambda\cong\mathbb{Z}_{p}[[T]]$ is a
two-dimensional regular local Noetherian ring with maximal ideal
$\mathfrak{m}=(p,T)$ and residue field $\mathbb{F}_{p}$. Its
$\mathfrak{m}$-adic topology coincides with the profinite topology.
\end{lemma}

\begin{proof}
This is Lang, \emph{Algebraic Number Theory}, Ch.~VII \S5. Regularity
follows from the embedding into $\mathbb{Z}_{p}[[T]]$ which is a
power-series ring over a regular local ring; dimension is $2$; the
maximal ideal is finitely generated by $(p,T)$; the topology
coincidence follows because both topologies are determined by
$\mathfrak{m}^{n}=(p,T)^{n}\supset p^{n}\Lambda+T^{n}\Lambda$.
\end{proof}

\subsection{Artin--Rees Lemma}
\label{v4-arith:w5-i:artin-rees}

\begin{lemma}[Artin--Rees for Noetherian filtered modules]
\label{v4-arith:w5-i:lem:artin-rees}
\ClaimStatusProvedHere \emph{(classical, Atiyah--Macdonald
Ch.~10 \S10.11)}. Let $R$ be a Noetherian ring, $\mathfrak{a}\subset R$
an ideal, $M$ a finitely generated $R$-module, and $N\subset M$ a
submodule. There exists an integer $k\geq 0$ such that for all
$n\geq k$,
\begin{equation}
\label{v4-arith:w5-i:eq:artin-rees}
\mathfrak{a}^{n}M\cap N\;=\;\mathfrak{a}^{n-k}(\mathfrak{a}^{k}M\cap N).
\end{equation}
\end{lemma}

\begin{proof}
Standard; the Rees algebra $\bigoplus_{n}\mathfrak{a}^{n}M$ is a
finitely generated module over the Rees algebra
$\bigoplus_{n}\mathfrak{a}^{n}$ (which is Noetherian by the
Hilbert basis theorem when $R$ is Noetherian), and the submodule
$\bigoplus_{n}(\mathfrak{a}^{n}M\cap N)$ is therefore finitely
generated; \eqref{v4-arith:w5-i:eq:artin-rees} follows by
generator-degree truncation. See Atiyah--Macdonald Prop.~10.9.
\end{proof}

\subsection{R1.formal.b via Artin--Rees}
\label{v4-arith:w5-i:R1-formal-b-closure}

We close (b.1) and (b.2) of
\Cref{v4-arith:bzn-sub:prop:R1-formal-b-scope}, at the level of
Iwasawa-lattice approximation.

\begin{proposition}[(b.2) holds for $\mathfrak{m}$-adically complete
ind-colimits]
\label{v4-arith:w5-i:prop:b2-closure}
\ClaimStatusProvedHere. Let $\mathcal{X}=\varinjlim_{i}\mathcal{X}_{i}$
be an ind-colimit of derived formal algebraic stacks with
$\mathcal{X}_{i}$ each finite-type and $\mathcal{X}$
$\mathfrak{m}$-adically complete along its formal locus. The natural
map
\[
\varinjlim_{i}\mathrm{Map}(B\widehat{\mathbb{Z}},\mathcal{X}_{i})\;\to\;\mathrm{Map}(B\widehat{\mathbb{Z}},\mathcal{X})
\]
is an equivalence in the pro-\'etale topology on the
$\mathfrak{m}$-adically-complete sector.
\end{proposition}

\begin{proof}
Each $B(\mathbb{Z}/p^{n}\mathbb{Z})$ is compact in the finite-type
sector, so $\mathrm{Map}(B(\mathbb{Z}/p^{n}\mathbb{Z}),-)$ commutes
with filtered colimits on the finite-type side (Lurie HTT~5.3.4.18).
$B\widehat{\mathbb{Z}}=\varprojlim_{n}B(\mathbb{Z}/p^{n}\mathbb{Z})$
is a cofiltered limit of compact objects. The key issue is whether
the $\mathrm{Map}$ functor commutes with this limit.

On the $\mathfrak{m}$-adically complete sector, the Artin--Rees
\Cref{v4-arith:w5-i:lem:artin-rees} applied to the ideal
$\mathfrak{m}=(p,T)\subset\Lambda$ and the ind-presentation
$\{\mathcal{X}_{i}\}$ gives: for each $n\geq 0$, there is
$k_{n}\geq 0$ such that the $(p,T)^{n}$-torsion in
$\mathcal{O}(\mathrm{Map}(B\widehat{\mathbb{Z}},\mathcal{X}))$
equals the image of
$\mathcal{O}(\mathrm{Map}(B(\mathbb{Z}/p^{k_{n}}\mathbb{Z}),\mathcal{X}_{i}))$
for some $i\gg 0$. Combined with the finite-type colimit exchange
for each compact $B(\mathbb{Z}/p^{k_{n}}\mathbb{Z})$, this yields
the colimit-commutation in each $\mathfrak{m}$-adic truncation;
taking the $\mathfrak{m}$-adic completion recovers the full
statement because the ind-colimit is
$\mathfrak{m}$-adically complete by hypothesis.
\end{proof}

\begin{corollary}[(b.1) holds on the
$\mathfrak{m}$-adically-complete sector]
\label{v4-arith:w5-i:cor:b1-closure}
\ClaimStatusProvedHere. On the $\mathfrak{m}$-adically-complete
sector of derived pro-\'etale stacks, the internal mapping stack
$\mathrm{Map}(B\widehat{\mathbb{Z}},-)$ commutes with filtered
colimits.
\end{corollary}

\begin{proof}
\Cref{v4-arith:w5-i:prop:b2-closure} establishes the specific
ind-colimit commutation; the general statement (b.1) is the case
of an arbitrary filtered diagram $\{\mathcal{Y}_{i}\}$ in the
$\mathfrak{m}$-adically complete sector. The same Artin--Rees
argument (applied to the diagram's structure rings
$\mathcal{O}(\mathrm{Map}(B(\mathbb{Z}/p^{k}\mathbb{Z}),\mathcal{Y}_{i}))$)
gives (b.1).
\end{proof}

\begin{theorem}[R1.formal.b closes on the completed sector]
\label{v4-arith:w5-i:thm:R1-formal-b-closes}
\ClaimStatusProvedHere. For ind-colimits of derived formal algebraic
stacks in the $\mathfrak{m}$-adically-complete sector, the R1.formal.b
statement
\[
L^{\mathrm{Ar}}\mathcal{X}\;\simeq\;\varinjlim_{i}L^{\mathrm{Ar}}\mathcal{X}_{i}
\]
holds.
\end{theorem}

\begin{proof}
Combine \Cref{v4-arith:w5-i:cor:b1-closure} (the general
filtered-colimit exchange) with
\Cref{v4-arith:w5-i:prop:b2-closure} (the specific presentation
exchange) via the identity
$L^{\mathrm{Ar}}\mathcal{X}=\mathrm{Map}(B\widehat{\mathbb{Z}},\mathcal{X})$.
The restriction to the $\mathfrak{m}$-adically-complete sector is
inherent to the Iwasawa-algebra setup and is the scope used in
Ch~32's R1.formal.b statement.
\end{proof}

\subsection{R1.formal.c.2 via Artin--Rees}
\label{v4-arith:w5-i:R1-formal-c-closure}

We close the colimit-commutation component (c.2) of
\Cref{v4-arith:bzn-sub:prop:R1-formal-c-scope}. The remaining
components (c.1) Arinkin--Gaitsgory extension and (c.3) flat-locus
matching with Fargues--Scholze are not closed by elementary
Artin--Rees arguments and remain open at their stated $5$--$10$-page
scope.

\begin{proposition}[(c.2) closes on the
$\mathfrak{m}$-adically-complete sector]
\label{v4-arith:w5-i:prop:c2-closure}
\ClaimStatusProvedHere. Let $\mathcal{X}=\varinjlim_{i}\mathcal{X}_{i}$
be an ind-colimit of derived formal algebraic stacks in the
$\mathfrak{m}$-adically-complete sector. The flat locus on the
arithmetic loop space commutes with the ind-colimit:
\[
L^{\mathrm{Ar}}\mathcal{X}^{\mathrm{flat}}\;\simeq\;\varinjlim_{i}L^{\mathrm{Ar}}\mathcal{X}_{i}^{\mathrm{flat}}.
\]
\end{proposition}

\begin{proof}
The flat locus is a closed subcategory in the $\mathrm{IndCoh}$
sense. By the Kan-extension property of closed subcategories under
filtered colimits (Lurie HA~4.2.4), the colimit
$\varinjlim_{i}L^{\mathrm{Ar}}\mathcal{X}_{i}^{\mathrm{flat}}$
embeds into $L^{\mathrm{Ar}}(\varinjlim_{i}\mathcal{X}_{i})^{\mathrm{flat}}$.
Conversely, the flat locus is preserved under the $(p,T)^{n}$-adic
truncations: an object is in the flat locus iff its
$(p,T)^{n}$-adic truncation is flat for every $n$, which reduces to
a finite-type compactness statement. By Artin--Rees
(\Cref{v4-arith:w5-i:lem:artin-rees}) applied to the flat-locus
submodule inside $\mathcal{O}(L^{\mathrm{Ar}}\mathcal{X})$, the
flat-locus filtration stabilises at some $k\geq 0$ relative to the
$\mathfrak{m}$-adic filtration; the stabilisation gives the reverse
embedding.
\end{proof}

\begin{remark}[(c.1) and (c.3) remain open]
\label{v4-arith:w5-i:rem:c1-c3-open}
Component (c.1) is the Arinkin--Gaitsgory extension of nilpotent
singular support to derived formal algebraic stacks; this is a
routine but non-trivial check in the Gaitsgory--Rozenblyum DAG Vol.~II
framework, estimated at $5$--$10$ pages in
\Cref{v4-arith:bzn-sub:rem:R1-formal-c-precise}. Component (c.3)
is the matching of the Arinkin--Gaitsgory flat locus with the
Fargues--Scholze $(\varphi,\Gamma)$-module flat locus, also a
$5$--$10$-page verification. Neither is closed by Artin--Rees alone;
both are marked open at inscription and not addressed further by
route~5 route~I.
\end{remark}

\subsection{Route~5 Route~I Formal-Noetherian Verdict}
\label{v4-arith:w5-i:formal-verdict}

\begin{theorem}[formal-Noetherian residual summary]
\label{v4-arith:w5-i:thm:formal-summary}
\ClaimStatusProvedHere. The formal-Noetherian residuals of
Ch~32 close as follows on the
$\mathfrak{m}$-adically-complete sector:
\begin{enumerate}
\item R1.formal.a: unconditional
(\Cref{v4-arith:bzn-sub:thm:R1-formal-a-closes}).
\item R1.formal.b: closed on the
$\mathfrak{m}$-adically-complete sector
(\Cref{v4-arith:w5-i:thm:R1-formal-b-closes}).
\item R1.formal.c.2: closed on the
$\mathfrak{m}$-adically-complete sector
(\Cref{v4-arith:w5-i:prop:c2-closure}).
\item R1.formal.c.1 and R1.formal.c.3: remain open, at stated
$5$--$10$-page scope each.
\end{enumerate}
\end{theorem}

\begin{proof}
(1) is cited; (2) and (3) are the content of the preceding
subsections; (4) is
\Cref{v4-arith:w5-i:rem:c1-c3-open}.
\end{proof}

\section{Summary}
\label{v4-arith:w5-i:summary}

The route~5 route~I slot closes three residuals identified in
Ch~44 and Ch~32:

\begin{enumerate}
\item $\mathrm{R}_{\mathrm{analytic}}$ \emph{density-of-zeros
component}: the sharp Riemann--von Mangoldt identity
\eqref{v4-arith:w5-i:eq:rvm} with $|S(T)|=O(\log T)$ is
unconditional
(\Cref{v4-arith:w5-i:thm:rvm},
\Cref{v4-arith:w5-i:thm:ST-bound},
\Cref{v4-arith:w5-i:cor:NT-sharp}). This closes the
density-of-zeros FRONTIER input; the positivity/linear-upper-growth
input remains open as the RH-adjacent half of
$\mathrm{R}_{\mathrm{analytic}}$.
\item \emph{Weil explicit formula at Gaussians}: the closed-form
evaluation \eqref{v4-arith:w5-i:eq:weil-gauss-closed} and
\eqref{v4-arith:w5-i:eq:weil-id-gauss} are inscribed; Gaussian
positivity \eqref{v4-arith:w5-i:eq:gauss-pos} holds trivially
(with equality) on the Gaussian convolution sub-class,
identifying the Gaussian class as too narrow to detect off-line
zeros (\Cref{v4-arith:w5-i:rem:gauss-pos-scope}). The explicit
closed form furnishes residual (H)'s computed reference
check-point.
\item \emph{R1.formal.b/c formal-Noetherian residuals}: closed on
the $\mathfrak{m}$-adically-complete sector via Artin--Rees
applied to the Iwasawa algebra
(\Cref{v4-arith:w5-i:thm:R1-formal-b-closes},
\Cref{v4-arith:w5-i:prop:c2-closure}). R1.formal.c.1 and
R1.formal.c.3 remain open at their stated $5$--$10$-page scope.
\end{enumerate}

The three closures are unconditional at their stated scope: (1) uses
only contour integration plus Stirling; (2) uses only Gaussian
self-duality plus Burnol--Cartier--Voros extension of Weil; (3) uses
only Artin--Rees on a Noetherian local ring. No RH content is
claimed; no H-P$^{\mathrm{Ar}}$ closure is claimed; the route~5
route~I contribution is the sharp density-of-zeros frontier, the
computed Weil-at-Gaussians reference, and the elementary Iwasawa
closure.

\bigskip

The route~5 route~I deliverable therefore tightens Ch~44's
three-axis classification of $\mathrm{R}_{\mathrm{analytic}}$ by
supplying the density-of-zeros axis in sharp unconditional form,
provides the explicit Gaussian Weil check-point at residual (H), and
closes the Artin--Rees-accessible components of the
formal-Noetherian residuals. The remaining RH-adjacent obstructions
(positivity of the regularised Hamburger moment functional, linear
upper-growth $\mu_{2n}^{1/(2n)}=O(n)$, Arinkin--Gaitsgory
derived-formal extension, Fargues--Scholze flat-locus matching)
remain open at the scope specified in Ch~41, Ch~44, and Ch~32.

\section{Bibliography}
\label{v4-arith:w5-i:bibliography}

Primary sources for this chapter:

\begin{description}
\item[Riemann 1859:] \"Uber die Anzahl der Primzahlen unter einer
gegebenen Gr\"o{\ss}e. Monatsber.~Berlin Akad.
\item[von Mangoldt 1895:] Zu Riemann's Abhandlung \"Uber die Anzahl
der Primzahlen unter einer gegebenen Gr\"o{\ss}e.
J.~Reine Angew.~Math.~114, 255--305.
\item[Selberg 1944:] On the remainder in the formula for $N(T)$,
the number of zeros of $\zeta(s)$ in the strip $0<t<T$.
Avh.~Norske Vid.~Akad.~Oslo, 1--27.
\item[Titchmarsh 1986:] \emph{The Theory of the Riemann Zeta-Function}
(2nd edn, rev.~Heath--Brown, Oxford).
\item[Weil 1952:] Sur les formules explicites de la th\'eorie des
nombres premiers. Comm.~S\'em.~Math.~Univ.~Lund, Tome suppl\'em.,
252--265.
\item[Tate 1950:] Fourier analysis in number fields and Hecke's
zeta-functions. PhD thesis, Princeton; in Cassels--Fr\"ohlich,
\emph{Algebraic Number Theory} (1967), Ch.~XV.
\item[Burnol 2000:] Sur certains espaces de Hilbert de fonctions
enti\`eres, li\'es \`a la transformation de Fourier.
C.~R.~Acad.~Sci.~Paris S\'er.~I 331, 919--924.
\item[Cartier--Voros 1990:] Une nouvelle interpr\'etation de la formule
des traces de Selberg. In \emph{The Grothendieck Festschrift} Vol.~II,
Progr.~Math.~87, 1--67.
\item[Atiyah--Macdonald 1969:] \emph{Introduction to Commutative
Algebra} (Addison-Wesley), Ch.~10.
\item[Bourbaki 1961:] \emph{Alg\`ebre Commutative}, Ch.~III \S3
(Masson).
\item[Lang 1994:] \emph{Algebraic Number Theory} (2nd edn, Springer),
Ch.~VII \S5.
\item[Edwards 1974:] \emph{Riemann's Zeta Function} (Academic Press).
\end{description}

These sources are disjoint from the route~4 audit catalog
(Akhiezer, Shohat--Tamarkin, Bombieri--Lagarias, Bost--Connes
references) in the Vol~IV sense: no route~4 primary derivation is
cited as a route~5-route-I verification, and vice versa.
